\section{Background and Related Work}
\label{sec:related}

\paragraph{Anomaly detection for IoT traffic}
Kitsune introduced a family of damped, incremental statistics computed per
packet over several time scales and several stream aggregations (host,
MAC--IP pair, channel, socket), together with an ensemble of small
autoencoders that learns benign behaviour online~\cite{mirsky2018kitsune}. It
was designed for resource-constrained network devices and evaluated on a
Raspberry~Pi. N-BaIoT applied the same 115 features to traffic captured from
nine commercial IoT devices before and after infection by the Mirai and
BASHLITE botnets, training one deep autoencoder per device and reporting that
all attacks were detected~\cite{meidan2018nbaiot}. Newer benchmarks such as
CICIoT2023 cover more devices and attack types~\cite{neto2023ciciot}.
\textsc{Sentinel} reuses the Kitsune feature definition and the N-BaIoT data,
but trains a single sequence autoencoder shared across devices, which is the
realistic choice for an appliance that cannot collect labelled benign traffic
for every new gadget in advance.

\paragraph{How such detectors are evaluated}
Sommer and Paxson argued that anomaly detection for networks is harder than
the machine-learning framing suggests: the cost of false positives is high,
the semantic gap between an anomaly and an attack is wide, and evaluation data
rarely resembles deployment~\cite{sommer2010closedworld}. Arp~et~al.\
catalogue recurring pitfalls in security machine learning, among them data
snooping, spurious correlations, inappropriate baselines and misleading
metrics under class imbalance~\cite{arp2022dosdonts}, and dataset audits have
found labelling and construction errors in widely used intrusion
datasets~\cite{engelen2021cicids}. Our evaluation is an application of that
advice to a complete system rather than to a single model; several of the
problems we report are instances of the pitfalls they describe.

\paragraph{LLMs for security operations}
Large language models are increasingly applied to security tasks, including
threat-intelligence extraction and mapping to ATT\&CK; a recent systematic
review surveys this literature~\cite{xu2025llmsec}, and surveys of ATT\&CK's
use in research describe the mapping problem in
detail~\cite{roy2023sok,alsada2023attack}. Agentic designs that chain
detection, classification and response have also appeared; AgenticCyber, for
instance, orchestrates cloud-hosted multimodal models over enterprise
telemetry~\cite{roy2025agenticcyber}. \textsc{Sentinel} differs in placing a
small quantised model on the appliance itself, which removes the need to send
traffic to a cloud service but constrains model size and speed. Most prior
evaluations of LLM ATT\&CK mapping use written threat reports as input; we
evaluate the model on the structured anomaly events a detector actually
produces.

\paragraph{Deception and automated blocking}
Honeypots are a long-standing source of attacker intelligence for IoT.
HoneyIoT adapts a high-interaction IoT honeypot's responses with
reinforcement learning to prolong attacker engagement~\cite{guan2023honeyiot}.
Heiding~et~al.\ used honeypot observations to evaluate login blocking for IoT
devices~\cite{heiding2020honeypot}, so feeding deception data into blocking is
not new. \textsc{Sentinel}'s contribution here is integration rather than
invention: honeypot sessions and detector alerts share one enforcement path
with one safety guard.

\paragraph{Kernel-level enforcement}
eBPF and XDP allow packets to be filtered in the kernel's fast path.
iKern embeds intrusion detection and prevention in eBPF~\cite{hadi2024ikern},
Zhang~et~al.\ run a neural network inside the kernel for real-time
prevention~\cite{zhang2024ebpf}, and eBPF-based DDoS mitigation has been
demonstrated on a Raspberry~Pi~4~\cite{tolay2025ebpf}. \textsc{Sentinel}
keeps all learning in user space and uses nftables for enforcement; it
contains an XDP drop path, but that path was not enabled in the configuration
evaluated here.

\paragraph{Home IoT context}
IoT Inspector collected labelled traffic from smart-home devices at scale and
found device behaviour to be comparatively regular~\cite{huang2020iotinspector},
which is the premise of anomaly detection on the home segment. Manufacturer
Usage Descriptions let vendors declare the traffic a device
needs~\cite{lear2019mud}; they are complementary to anomaly detection, but
require manufacturer participation that most consumer devices lack.
